\documentclass[11pt,a4paper]{article}

\usepackage[utf8]{inputenc}
\usepackage[T1]{fontenc}
\usepackage[spanish,es-tabla]{babel}
\usepackage{amsmath,amssymb,amsthm,mathtools}
\usepackage{booktabs}
\usepackage{enumitem}
\usepackage{geometry}
\usepackage{setspace}
\usepackage{hyperref}
\usepackage{microtype}

\geometry{margin=1in}
\onehalfspacing

\hypersetup{
  pdftitle={Nota formal: Feigenbaum y estructura de conjunciones},
  pdfauthor={Johel Padilla-Villanueva}
}

\theoremstyle{definition}
\newtheorem{definition}{Definici\'on}[section]
\newtheorem{remark}{Observaci\'on}[section]
\newtheorem{example}{Ejemplo}[section]
\theoremstyle{plain}
\newtheorem{proposition}{Proposici\'on}[section]
\newtheorem{lemma}{Lema}[section]
\newtheorem{theorem}{Teorema}[section]
\newtheorem{corollary}{Corolario}[section]
\newtheorem{conjecture}{Conjetura}[section]
\theoremstyle{remark}
\newtheorem{scholium}{Scholium}[section]

\newcommand{\Res}{\mathrm{Res}}
\newcommand{\RECD}{\mathrm{RECD}}
\newcommand{\TC}{\mathrm{TC}}
\newcommand{\excess}{\mathrm{excess3}}
\newcommand{\Rpair}{\mathcal{R}_{\mathrm{pair}}}
\newcommand{\Jwin}{\mathcal{J}}
\newcommand{\ts}{\tau_s}

\title{\textbf{Nota formal}\\[0.4em]
\large Conjunciones ordinales a lo largo de la\\
ruta de Feigenbaum\\[0.6em]
\normalsize (Punto 4 del roadmap; asume Notas 1--3)}

\author{
Johel Padilla-Villanueva\\
\textit{Department of Environmental Health, University of Puerto Rico}\\
\texttt{johelpadilla@gmail.com}
}

\date{2026-08-01 \quad$\cdot$\quad v0.2 (barrido denso $r$)}

\begin{document}
\maketitle

\begin{abstract}
Esta nota caracteriza \emph{qu\'e cambia} en la estructura de
conjunciones ordinales cuando un sistema multivariante atraviesa la
cascada de doblamiento de periodo hacia el caos, en lugar de limitarse a
afirmar que ``hay conexi\'on con Feigenbaum''. Se fija el escenario de
mapas unimodales acoplados, se definen observables de profundidad
(abundancia vs share), se enuncian predicciones estratificadas
(pre-acumulaci\'on, cerca de $r_\infty$, caos desarrollado, ventanas
peri\'odicas), y se separa lo demostrado, lo conjeturado y lo solo
evidenciado sint\'eticamente. El contraste con EWS cl\'asicos
(varianza, AC1) aclara qu\'e ve Systemic Tau / RECD de distinto:
reorganizaci\'on \emph{relacional de profundidad}, no solo
ralentizaci\'on univariante.
\end{abstract}

\tableofcontents
\vspace{1em}

\section{Canon y pregunta}
\label{sec:q}

\begin{scholium}
Notas 1--3: claims/estimadores/reloj; L3 fuerte $:=$ $\Res_{\mathrm{pair}}$;
$\excess$ proxy; $\alpha$ motor admisible; abundancia $\ne$ share $f_3$;
M-I-1.
\end{scholium}

\textbf{Pregunta del punto~4:}
¿Qu\'e cambia \emph{estructuralmente} en las conjunciones
$(\Phi_1,\Phi_2,\Res_{\mathrm{pair}})$ en cada etapa de la cascada, y
qu\'e de eso no capturan las medidas cl\'asicas?

\section{Escenario din\'amico}
\label{sec:scene}

\begin{definition}[Familia log\'istica y acumulaci\'on]
\label{def:logistic}
$x\mapsto r x(1-x)$ en $[0,1]$, con cascada de doblamientos que
acumula en $r_\infty\approx 3.56995$ y raz\'on asint\'otica
$\delta\approx 4.6692016091$.
\end{definition}

\begin{definition}[Mapas acoplados de trabajo]
\label{def:coupled}
Para $d\ge 3$ y acoplamiento difusivo d\'ebil $\varepsilon>0$,
\begin{equation}
x_{n+1}^{(i)}
=(1-\varepsilon)\,f_r\bigl(x_n^{(i)}\bigr)
+\frac{\varepsilon}{d-1}\sum_{j\ne i}f_r\bigl(x_n^{(j)}\bigr)
+\eta_n^{(i)},
\end{equation}
con $f_r$ unimodal de la clase Feigenbaum (cr\'itico cuadr\'atico) y
$\eta$ ruido de observaci\'on peque\~no. Defaults de validaci\'on del
paper de fundamentos: $d=4$, $\varepsilon=0.05$, $\sigma\approx 0.003$.
\end{definition}

\begin{definition}[Estaciones de la cascada]
\label{def:stations}
\begin{enumerate}[label=(\roman*)]
  \item \textbf{S$_k$:} r\'egimen de periodo $2^k$ (atractor peri\'odico
  estable de periodo $2^k$), $r\in(r_k,r_{k+1})$.
  \item \textbf{S$_\infty^-$:} pre-acumulaci\'on, $r\uparrow r_\infty$
  con $k$ grande.
  \item \textbf{S$_\infty$:} vecindad de la acumulaci\'on.
  \item \textbf{S$_{\mathrm{ch}}$:} caos desarrollado (p.ej.\ $r=3.85$
  en la validaci\'on).
  \item \textbf{S$_{\mathrm{win}}$:} ventanas peri\'odicas embebidas en
  caos.
\end{enumerate}
\end{definition}

\begin{definition}[Proceso de s\'imbolos y observables]
\label{def:obs}
$S_t$ Bandt--Pompe como en Nota~1. Observables de ventana:
\begin{align}
A_1&:=\mathbb{E}[\Phi_1],&
A_2&:=\mathbb{E}[\Phi_2],&
A_3&:=\mathbb{E}[\Res_{\mathrm{pair}}]
\text{ o }\mathbb{E}[\excess],\\
\rho_3&:=\mathbb{P}(\excess>\theta_{\mathrm{hi}}),&
\overline{\ts}&:=\mathbb{E}[|\ts|],&
f_3&:=\mathbb{E}\bigl[\text{share L3 bajo }\alpha(\lambda)\bigr].
\end{align}
$A_3,\rho_3$ son \textbf{abundancia}; $f_3$ es \textbf{share}
(Nota~3).
\end{definition}

\section{Qu\'e cambia en cada estaci\'on (predicciones)}
\label{sec:pred}

\subsection{Estructura esperada por estaci\'on}

\begin{center}
\begin{tabular}{@{}p{0.12\textwidth}p{0.28\textwidth}p{0.28\textwidth}p{0.22\textwidth}@{}}
\toprule
Estaci\'on & Organizaci\'on ordinal t\'ipica & Profundidad dominante & EWS cl\'asicos \\
\midrule
S$_k$ (bajo $k$) &
Patrones limpios, alta coincidencia y relaciones estables &
$\Phi_1,\Phi_2$ altos; $\Res_{\mathrm{pair}}$ bajo &
Var/AC1 bajos tras transitorio \\
S$_\infty^-$ &
Periodos largos; frontera de predictibilidad local &
$\Phi_2$ puede sostenerse; L3 empieza a crecer &
CSD: var/AC1 $\uparrow$ al acercarse a $r_{k+1}$ \\
S$_\infty$ &
Mezcla de escalas; colapso de predictibilidad local &
Reorganizaci\'on; banda $|\ts|$ inestable &
CSD m\'aximo local; ambig\"uedad multivariante \\
S$_{\mathrm{ch}}$ &
Mezcla fuerte; joint no-pair &
$A_3,\rho_3,f_3$ elevados (pred.\ central) &
Var alta; no ve profundidad ordinal \\
S$_{\mathrm{win}}$ &
Retorno parcial a S$_k$-like &
Ca\'ida parcial de $A_3$ &
Var puede caer en ventana \\
\bottomrule
\end{tabular}
\end{center}

\begin{conjecture}[Monoton\'ia gruesa pre$\to$chaos de abundancia L3]
\label{conj:mono}
En la clase Def.~\ref{def:coupled} con $\varepsilon$ peque\~no y $d\ge 3$,
\begin{equation}
\mathbb{E}\bigl[A_3\mid\mathrm{S}_{\mathrm{ch}}\bigr]
>
\mathbb{E}\bigl[A_3\mid\mathrm{S}_{k}\bigr]
\end{equation}
para $k$ acotado (p.ej.\ periodo 2 o 4) y $A_3\in\{\mathbb{E}\excess,
\mathbb{E}\Res_{\mathrm{pair}}\}$. Compatible con la validaci\'on
sint\'etica ($r=3.30$ vs $r=3.85$: mean excess3 $1.76\to 1.89$, highL3
$0.50\to 0.90$, $f_3$ $0.61\to 0.87$ bajo $\lambda(r)$).
\end{conjecture}

\begin{conjecture}[No trivialidad respecto de $|\ts|$ y $\Phi_1$]
\label{conj:nontriv}
El incremento $A_3(\mathrm{S}_{\mathrm{ch}})-A_3(\mathrm{S}_{k})$
\textbf{no} es una funci\'on mon\'otona fija de
$\Delta\overline{|\ts|}$ ni de $\Delta A_1$. Existen realizaciones /
acoplamientos donde $|\ts|$ cae y $A_3$ sube (desincronizaci\'on con
surplus de orden 3), y el rec\'iproco en drivers comunes (Nota~2).
\end{conjecture}

\begin{conjecture}[Acumulaci\'on $r_\infty$]
\label{conj:rinf}
En una vecindad de $r_\infty$, la \emph{variabilidad temporal} de
$\Res_{\mathrm{pair}}(t)$ y de $\ts(t)$ se maximiza localmente (r\'egimen
de reorganizaci\'on), mientras que el valor medio de $A_3$ puede a\'un
estar por debajo del plateau de caos desarrollado. Es decir: cerca de
la acumulaci\'on domina la \textbf{inestabilidad de la geometr\'ia
relacional}; en S$_{\mathrm{ch}}$ domina la \textbf{masa media de L3}.
\end{conjecture}

\begin{conjecture}[Doblamiento $2^k\to 2^{k+1}$]
\label{conj:doubling}
En cada doblamiento, el grafo de relaciones persistentes ($\Phi_2$) se
reconfigura (cambian c\'odigos dominantes $R^{ij}$) en una escala de
tiempo ligada al nuevo periodo, mientras $\Phi_1$ puede oscilar sin
tendencia mon\'otona universal. El residuo $\Res_{\mathrm{pair}}$ no
necesita saltar en cada $r_k$; su crecimiento sistem\'atico se concentra
al acumularse la cascada y al entrar en caos.
\end{conjecture}

\section{Qu\'e ve el paradigma de distinto a EWS cl\'asicos}
\label{sec:vs_ews}

\begin{definition}[Familia EWS cl\'asica]
\label{def:ews}
Varianza local, autocorrelaci\'on lag-1, a veces skewness/recovery rate:
firmas de \emph{critical slowing down} (CSD) univariante o
semi-univariante.
\end{definition}

\begin{theorem}[Separaci\'on de objetos de medida]
\label{thm:sep}
\begin{enumerate}[label=(\roman*)]
  \item EWS cl\'asicos miden propiedades de la \emph{ley marginal
  temporal} (o de un agregado univariante) cerca de una bifurcaci\'on
  local.
  \item $\ts$ mide \emph{concordancia ordinal entre canales}
  (observable de orden $\le 2$ entre series).
  \item $(\Phi_1,\Phi_2,\Res_{\mathrm{pair}})$ miden
  \emph{profundidad de conjunción} en el joint simb\'olico.
  \item Por tanto, un sistema puede exhibir CSD fuerte con L3 bajo
  (bifurcaci\'on univariante sin sinergia de orden 3), o L3 alto sin CSD
  cl\'asico claro (caos desarrollado ya lejos del punto cr\'itico local,
  o transici\'on relacional sin slowing univariante dominante).
\end{enumerate}
\end{theorem}

\begin{proof}[Esbozo]
(i)--(iii) por definici\'on de los funcionales. (iv) sigue de la
independencia formal de los argumentos: los EWS no ven $\Sigma^d$; L3
no requiere un autovalor de Jacobiano cerca de 1.
\end{proof}

\begin{proposition}[Rol dual $\ts$ / RECD en la cascada]
\label{prop:dual}
\begin{enumerate}[label=(\roman*)]
  \item $\ts$ responde: \emph{¿se reorganiza el acoplamiento ordinal
  mutuo?}
  \item RECD / $(A_\ell,f_\ell)$ responde: \emph{¿a qu\'e profundidad
  se expresa esa reorganizaci\'on?}
  \item En S$_{\mathrm{ch}}$, la predicci\'on conjunta es: posible ca\'ida
  o volatilidad de $|\ts|$ \emph{y} subida de $A_3$ (no
  contradictorias: orden $\le 2$ vs residuo).
\end{enumerate}
\end{proposition}

\begin{scholium}
Systemic Tau \textbf{no} es ``otro indicador de variance''. Detecta
estructura relacional; el RECD anidado a\~nade el eje de
\emph{profundidad}. Eso es lo estructuralmente distinto en la cascada.
\end{scholium}

\section{Papel de $\delta$ y de $\alpha(\lambda)$}
\label{sec:delta}

\begin{proposition}[Dos usos de $\delta$]
\label{prop:delta_uses}
\begin{enumerate}[label=(\roman*)]
  \item \textbf{Ancla de escala en el cue:}
  $\kappa_\delta=1+\gamma\log\delta$ dentro de $\lambda$ (Nota~3) ---
  continuidad con renormalización, no medici\'on de $\delta$ desde
  series cortas.
  \item \textbf{Teorema de reducci\'on temporal (rama legacy / Tomus II):}
  contracci\'on $\delta^{-k}$ en construcciones de paso de tiempo a lo
  largo de la cascada --- objeto distinto del anidado
  $\sum\alpha_\ell\Phi_\ell$.
\end{enumerate}
No deben fusionarse sin scholium (M-II-4 / pureza Tomus II).
\end{proposition}

\begin{proposition}[Abundancia vs share en la cascada]
\label{prop:ab_share}
Con $\lambda=\lambda(r)$ mon\'otono que cumple el design goal:
\begin{enumerate}[label=(\roman*)]
  \item $A_3(r)$ es predicci\'on de \emph{objeto} (Notas 1--2);
  \item $f_3(r)$ mezcla $A_3$ con $\alpha_3(\lambda(r))$ (Nota~3);
  \item un aumento de $f_3$ con $\alpha$ admisible y $A_3$ plano a\'un
  es solo reponderaci\'on, no nueva evidencia de sinergia.
\end{enumerate}
La validaci\'on sint\'etica del paper reporta ambos; la pretensi\'on
te\'orica fuerte recae en (i).
\end{proposition}

\section{Evidencia sint\'etica existente (estatus)}
\label{sec:evidence}

\subsection{Contraste binario pre$\to$caos (paper de fundamentos)}

\begin{center}
\begin{tabular}{@{}llll@{}}
\toprule
M\'etrica & Pre $r=3.30$ & Chaos $r=3.85$ & Estatus \\
\midrule
mean excess3 & $1.760\pm 0.004$ & $1.892\pm 0.006$ & Evidencia (no thm) \\
highL3 ($>1.75$) & $0.502$ & $0.904$ & Evidencia \\
$f_3$ (con $\lambda(r)$) & $0.613$ & $0.870$ & Evidencia + motor \\
$\lambda_{\mathrm{emp}}$ sin calibrar & $\approx 0.181$ & $\approx 0.108$ &
Fragilidad del cue $|\ts|$ \\
\bottomrule
\end{tabular}
\end{center}

\subsection{Barrido denso de $r$ a lo largo de la cascada (v0.2)}
\label{sec:sweep}

Protocolo: $d=4$, $\varepsilon=0.05$, $\sigma=0.003$, $T=2200$ post-burn,
$m=3$, $w=13$, $N_{\mathrm{real}}=6$ por punto, rejilla de $32$ valores de
$r\in[3.10,3.95]$ densificada cerca de $r_\infty\approx 3.56995$.
Software: \texttt{nested-recd} (excess3 proxy; $\Res_{\mathrm{pair}}$ con
stride~$8$). Script: \texttt{scripts/cascade\_r\_sweep.py}.
Figuras estables: \texttt{figures/cascade\_A3\_vs\_r.png},
\texttt{cascade\_variability\_tau.png}, \texttt{cascade\_ews\_vs\_A3.png},
\texttt{cascade\_phi\_f3.png}.

\begin{center}
\small
\begin{tabular}{@{}lcccccc@{}}
\toprule
Estaci\'on (aprox.) & $r$ & excess3 & highL3 & $\Res_{\mathrm{pair}}$ &
$|\ts|$ & var \\
\midrule
S$_k$ (p2) & $3.20$ & $1.754$ & $0.47$ & $0.199$ & $0.54$ & $0.021$ \\
S$_k$ (p2) & $3.30$ & $1.758$ & $0.47$ & $0.187$ & $0.17$ & $0.025$ \\
S$_k$ (p4) & $3.45$ & $1.964$ & $0.85$ & $0.039$ & $0.43$ & $0.040$ \\
S$_\infty^-$ & $3.564$ & $2.373$ & $0.99$ & $\approx 0$ & $0.39$ & $0.046$ \\
S$_\infty$ & $3.570$ & $2.343$ & $0.99$ & $\approx 0$ & $0.20$ & $0.045$ \\
S$_{\mathrm{ch}}$ & $3.70$ & $1.945$ & $0.89$ & $0.007$ & $0.07$ & $0.050$ \\
S$_{\mathrm{ch}}$ & $3.85$ & $1.863$ & $0.84$ & $0.008$ & $0.12$ & $0.056$ \\
S$_{\mathrm{ch}}$ & $3.95$ & $1.899$ & $0.88$ & $0.007$ & $0.12$ & $0.080$ \\
\bottomrule
\end{tabular}
\end{center}

\paragraph{Lectura operativa de las conjeturas.}
\begin{enumerate}[label=(\roman*)]
  \item \textbf{Conj.~\ref{conj:mono} (proxy excess3 / highL3):}
  sostenida en el contraste grueso $r=3.30\to 3.85$
  (excess3 $1.76\to 1.86$; highL3 $0.47\to 0.84$). El \emph{m\'aximo}
  de mean excess3 en la rejilla no est\'a en caos profundo sino cerca de
  S$_\infty^-$ ($r\approx 3.566$, excess3 $\approx 2.38$). Es decir:
  monotonia gruesa pre$\to$chaos s\'i; plateau de caos desarrollado
  \textbf{no} es el pico de abundancia media del proxy.
  \item \textbf{Conj.~\ref{conj:mono} (fuerte $\Res_{\mathrm{pair}}$):}
  con el protocolo del barrido ($w{=}13$) el residual parece m\'aximo
  en S$_k$ de periodo bajo y $\approx 0$ en caos --- interpretaci\'on
  \emph{provisional} del v0.2. El diagn\'ostico v0.3
  (\S\ref{sec:resdiag}) muestra que eso es en gran parte
  \textbf{sesgo de ventana}: con $w{=}104$ o bloque largo,
  $\Res(r{=}3.85)\gg\Res(r{=}3.20)$ y el residual en $r_\infty$
  permanece bajo. excess3 y $\Res_{\mathrm{pair}}$ siguen sin ser
  intercambiables (trayectorias y sensibilidad a $w$ distintas).
  \item \textbf{Conj.~\ref{conj:nontriv}:} sostenida.
  $\mathrm{corr}_r(\mathrm{excess3},|\ts|)\approx 0.27$,
  $\mathrm{corr}_r(\mathrm{excess3},\mathrm{var})\approx 0.21$;
  en $3.30\to 3.85$, $\Delta$excess3$>0$ mientras $\Delta|\ts|$ no es
  un rescaling mon\'otono del mismo signo a lo largo de toda la cascada.
  \item \textbf{Conj.~\ref{conj:rinf}:} parcialmente matizada.
  El pico de \emph{abundancia media} de excess3 est\'a en
  S$_\infty^-$/S$_\infty$; la $\mathrm{std}_t(\mathrm{excess3})$ m\'axima
  en la rejilla aparece algo antes ($r\approx 3.52$). El caos
  desarrollado retiene highL3 alto pero mean excess3 intermedio.
  \item \textbf{EWS cl\'asicos:} la varianza crece de forma m\'as
  mon\'otona hacia $r$ alto; su forma normalizada \textbf{no} replica la
  de excess3 (figura \texttt{cascade\_ews\_vs\_A3.png}).
\end{enumerate}

\begin{scholium}
Los n\'umeros binarios del paper \textbf{falsan} la nula gruesa ``L3 no
cambia pre$\to$chaos'' en excess3/highL3. El barrido denso \textbf{refina}
Conjs.~\ref{conj:mono}--\ref{conj:rinf}: el pico de abundancia del proxy
est\'a cerca de la acumulaci\'on, no en caos profundo.
El diagn\'ostico de $\Res_{\mathrm{pair}}$ (\S\ref{sec:resdiag})
\textbf{corrige} la lectura del residual fuerte: el colapso en caos
con $w{=}13$ era sesgo; con $w$ adecuado la monotonia gruesa se
restaura, mientras $r_\infty$ retiene residual estructuralmente bajo.
Nada de esto es teorema en la clase Feigenbaum acoplada ni en datos
emp\'iricos.
\end{scholium}

\begin{scholium}[Canon emp\'irico a\~nadido]
\label{sch:emp}
\begin{enumerate}[label=\textbf{E\arabic*.}]
  \item Reportar siempre $A_3$ de abundancia (excess3, highL3) y, cuando
  sea computable, $\Res_{\mathrm{pair}}$ por separado.
  \item No identificar el m\'aximo de excess3 con ``caos m\'as
  desarrollado'': en este ensamble el m\'aximo es pre/peri-acumulaci\'on.
  \item Contraste cl\'asico EWS vs A3 es obligatorio en cualquier
  validaci\'on de cascada (objetos distintos; Thm.~\ref{thm:sep}).
  \item No interpretar $\Res_{\mathrm{pair}}(w{=}13)\approx 0$ en caos
  como ausencia de L3 fuerte: diagnosticar con $w$ largo y bloque
  completo (E4--E6 abajo).
\end{enumerate}
\end{scholium}

\subsection{Diagn\'ostico del colapso de $\Res_{\mathrm{pair}}$ (v0.3)}
\label{sec:resdiag}

Protocolo (script \texttt{scripts/res\_pair\_diagnostics.py}):
estaciones $r\in\{3.20,3.30,3.45,3.57,3.70,3.85\}$,
$w\in\{13,26,52,104\}$, $N_{\mathrm{real}}=4$, $T=1800$,
m\'etricas por ventana: $\Res_{\mathrm{pair}}$, sparsity
($\#$ joints \'unicos$/w$), alfabeto, $H$, $\mathrm{TC}-\mathrm{TC}(P^{(2)})$,
desviaci\'on final del IPF, y residual de un bloque largo
($\sim 800$ s\'imbolos).

\begin{center}
\small
\begin{tabular}{@{}lcccccc@{}}
\toprule
$r$ & est. & $w{=}13$ & $w{=}26$ & $w{=}52$ & $w{=}104$ & full \\
\midrule
$3.20$ & S$_k$ p2 & $0.192$ & $0.317$ & $0.184$ & $0.080$ & $0.009$ \\
$3.30$ & S$_k$ p2 & $0.181$ & $0.303$ & $0.183$ & $0.080$ & $0.009$ \\
$3.45$ & S$_k$ p4 & $0.097$ & $0.200$ & $0.151$ & $0.069$ & $0.005$ \\
$3.57$ & S$_\infty$ & $0.001$ & $0.005$ & $0.008$ & $0.012$ & $0.001$ \\
$3.70$ & S$_{\mathrm{ch}}$ & $0.004$ & $0.016$ & $0.041$ & $0.062$ & $0.005$ \\
$3.85$ & S$_{\mathrm{ch}}$ & $0.009$ & $0.150$ & $0.940$ & $1.259$ & $0.454$ \\
\bottomrule
\end{tabular}
\end{center}

\paragraph{Bater\'ia de hip\'otesis (veredictos operativos).}
\begin{enumerate}[label=\textbf{H\arabic*.}]
  \item \textbf{Sesgo de ventana (parcial, decisivo en caos).}
  En $r=3.85$, $\Res$ pasa de $\approx 0.009$ ($w{=}13$) a
  $\approx 1.26$ ($w{=}104$) y $\approx 0.45$ en bloque completo.
  El ``colapso en caos'' del barrido denso con $w{=}13$ era en gran
  parte \emph{falso negativo del estimador}. En $r_\infty$, el lift
  es d\'ebil ($0.001\to 0.012$): all\'i el residual permanece bajo.
  \item \textbf{Sparsity alfabetica (no como se conjetur\'o).}
  En $w{=}13$, sparsity en S$_k$ p2 ($\approx 0.87$) es \emph{mayor}
  que en $r_\infty$ ($\approx 0.41$). Cerca de la acumulaci\'on el
  joint se \emph{contrae} (pocos patrones, ciclos de alto periodo
  casi sincronizados), no se expande. La correlaci\'on
  sparsity--$\Res$ es positiva ($\approx +0.47$), no negativa.
  \item \textbf{Fallo de IPF (rechazado).}
  $\mathrm{max\_dev}$ final del IPF $\ll 10^{-3}$ en todas las
  estaciones; el colapso no es no-convergencia num\'erica.
  \item \textbf{Pico de periodo-2 en $w{=}13$ es inflaci\'on de
  ventana corta.} El residual asint\'otico (bloque $\sim 800$) en
  p2 es $\approx 0.009$, del orden del de $r_\infty$. El ``pico''
  reportado en el barrido denso era sesgo de $w$ peque\~na sobre
  tablas joint con pocas repeticiones.
  \item \textbf{Conj.~\ref{conj:mono} fuerte (revisada).}
  Con $w\ge 52$ o bloque largo: $\Res(r{=}3.85)\gg\Res(r{=}3.20)$.
  La monotonia gruesa pre$\to$chaos \emph{s\'i} se sostiene para
  $\Res_{\mathrm{pair}}$ bajo estimaci\'on adecuada; falla solo
  bajo el protocolo $w{=}13$ del barrido inicial.
\end{enumerate}

\begin{scholium}[Canon diagn\'ostico]
\label{sch:resdiag}
\begin{enumerate}[label=\textbf{E\arabic*.}]
  \setcounter{enumi}{3}
  \item Todo claim de ``$\Res_{\mathrm{pair}}\approx 0$'' debe
  acompa\~narse de $w$ (o bloque) y de un chequeo de sensibilidad
  en $w$.
  \item Distinguir tres reg\'imenes del residual fuerte en esta
  clase: (i)~caos desarrollado = residual alto si $w$ basta;
  (ii)~$r_\infty$ = residual estructuralmente bajo (soporte
  conjunto reducido y casi pairwise-explicable);
  (iii)~periodo bajo + $w$ corto = inflaci\'on del residual.
  \item excess3 (proxy) y $\Res_{\mathrm{pair}}$ (fuerte) siguen
  trayectorias distintas en $r$ \emph{y} en $w$; no sustituir uno
  por el otro en claims (refuerzo de E1).
\end{enumerate}
\end{scholium}

Artefactos: \texttt{results/res\_pair\_diag/}, figuras
\texttt{figures/res\_pair\_diag\_*.png}.

\section{Programa de teoremas (punto 4 $\to$ agenda)}
\label{sec:program}

\begin{enumerate}[label=\textbf{T4.\arabic*.}]
  \item Cotas de $\Res_{\mathrm{pair}}$ para productos de medidas
  invariantes de mapas unimodales acoplados d\'ebilmente (l\'imite
  $\varepsilon\to 0$ vs $\varepsilon>0$ fijo).
  \item Continuidad/medida de $A_3(r)$ fuera de ventanas peri\'odicas.
  \item Teorema de no-reducci\'on: existencia de familias donde
  $\partial_r A_3$ no es funci\'on de $\partial_r\mathrm{Var}$ ni de
  $\partial_r\overline{|\ts|}$.
  \item Caracterizaci\'on de la varianza de $\Res_{\mathrm{pair}}(t)$ en
  $r_\infty$ (Conj.~\ref{conj:rinf}).
\end{enumerate}

\section{Respuesta al criterio de \'exito}
\label{sec:ans}

\begin{enumerate}[label=\textbf{Q\arabic*.}, leftmargin=*]
  \item \textbf{¿Qu\'e cambia estructuralmente cerca de una transici\'on
  cr\'itica en esta ruta?}\\
  La geometr\'ia relacional del joint simb\'olico: de describible por
  pares/persistencia (S$_k$) a masa creciente de residuo fuera de
  $\Rpair$ (S$_{\mathrm{ch}}$), con zona de m\'axima
  \emph{inestabilidad} de esa geometr\'ia cerca de $r_\infty$
  (Conj.~\ref{conj:rinf}).

  \item \textbf{¿Aparece el Nivel~3 de forma caracter\'istica cerca de
  la acumulaci\'on?}\\
  Hip\'otesis: s\'i como \emph{volatilidad} de L3 y de $\ts$; el
  \emph{plateau} de abundancia media es m\'as propio del caos
  desarrollado que del punto $r_\infty$ aislado.

  \item \textbf{¿Qu\'e ve Systemic Tau de distinto?}\\
  Eje relacional ($\ts$) + eje de profundidad
  $(\Phi_1,\Phi_2,\Res_{\mathrm{pair}})$, duales; EWS cl\'asicos no
  tienen eje de profundidad ordinal multivariante
  (Thm.~\ref{thm:sep}).
\end{enumerate}

\section{Canon a\~nadido (C23--C28)}
\label{sec:canon}

\begin{enumerate}[label=\textbf{C\arabic*.}]
  \setcounter{enumi}{22}
  \item Estaciones S$_k$, S$_\infty^-$, S$_\infty$, S$_{\mathrm{ch}}$,
  S$_{\mathrm{win}}$ para hablar de la cascada.
  \item Predicci\'on central: $A_3(\mathrm{S}_{\mathrm{ch}})>A_3(\mathrm{S}_k)$
  en la clase acoplada d\'ebil (conjetura + evidencia sint\'etica).
  \item No-trivialidad: L3 no es rescaling de $|\ts|$ ni de $\Phi_1$.
  \item Separaci\'on EWS cl\'asicos / $\ts$ / profundidad RECD
  (Thm.~\ref{thm:sep}).
  \item $\delta$ tiene dos usos (ancla de cue vs reducci\'on legacy);
  no fusionar.
  \item Reportar abundancia y share por separado a lo largo de $r$.
\end{enumerate}

\vspace{1.5em}
\noindent\rule{\textwidth}{0.4pt}\\[0.5em]
\begin{center}
\em Fin Nota formal punto 4 v0.3 (barrido denso + diagn\'ostico
$\Res_{\mathrm{pair}}$).\\
Artefactos: \texttt{results/cascade\_sweep/},
\texttt{results/res\_pair\_diag/}, figuras
\texttt{cascade\_*.png}, \texttt{res\_pair\_diag\_*.png};
software \texttt{nested-recd}.
\end{center}

\begin{thebibliography}{99}
\bibitem{Padilla2026framework}
J.~Padilla-Villanueva. Systemic Tau and the RECD Framework. 2026.
\bibitem{Feigenbaum1978}
M.~J.~Feigenbaum. Quantitative universality for a class of nonlinear
transformations. \emph{J.\ Stat.\ Phys.} 19, 25--52 (1978).
\bibitem{Kaneko1992}
K.~Kaneko. Overview of coupled map lattices.
\emph{Chaos} 2, 279--282 (1992).
\end{thebibliography}

\end{document}
